\begin{table*}[t]
    \centering
    \caption{Draft analysis roadmap. The goal is to keep each figure tied to one interpretable scientific claim rather than touring the whole dashboard.}
    \label{tab:analysis-roadmap}
    \small
    \begin{tabularx}{\textwidth}{p{0.09\textwidth}p{0.25\textwidth}p{0.31\textwidth}X}
        \toprule
        RQ & Scientific question & Figure or table & Interpretation target \\
        \midrule
        RQ0 & Is the actually used budget observable? & \tabref{tab:budget-observability}; \tabref{tab:used-budget-by-requested}; \figref{fig:used-reasoning-budget-latex} & Separate tokenizer-counted emitted reasoning length from requested budget before making cost claims. \\
        RQ1 & What is the baseline behavior across final, chain, and trace metrics? & \figref{fig:overall-metrics-latex}; \tabref{tab:overall} & Check whether final-answer accuracy hides weak reasoning traces. \\
        RQ2 & How does performance scale with requested budget? & \figref{fig:budget-scaling-latex}; \tabref{tab:budget} & Estimate budget sensitivity and compare model-specific gains. \\
        RQ3 & Where does the budget saturate? & \tabref{tab:marginal-budget} & Identify the first budget step where marginal gains become negligible. \\
        RQ4 & At 16,384 tokens, what remains hard? & \figref{fig:depth-context-latex}; \tabref{tab:high-budget-16384} & Study depth and context after giving the model maximal requested budget. \\
        RQ5 & Are the remaining errors structural? & \figref{fig:error-taxonomy-latex}; \tabref{tab:errors} & Distinguish wrong final decisions from wrong paths and invalid traces. \\
        \bottomrule
    \end{tabularx}
\end{table*}
